\section{Conclusion}
In this paper, we present Bind-Your-Avatar, an MM-DiT-based framework for multi-talking-character video generation in a shared scene. To address the challenge of audio-to-character correspondence, we propose a fine-grained embedding router that binds “who” and “speak what” together. We introduce two 3D-mask router implementations with geometry-based losses and mask refinement. To support this task, we construct the first dedicated dataset and release an open-source data processing pipeline. Extensive experiments on MTCC and our Bind-Your-Avatar-Benchmark demonstrate the effectiveness of our approach.
\subsubsection{Limitations}
Our model lacks explicit human motion modeling and is trained on limited audio-driven video data, which seldom include multi-person interactions. This results in suboptimal hand gestures and less realistic character interactions. Additionally, the full-attention architecture of MM-DiT leads to high inference costs, limiting real-time applicability. Future work will focus on improving motion fidelity and reducing computational overhead for practical deployment.
\appendix
\section{Intra-Denoise Router Architecture}
The network architecture of Intra-Denoise Router, grounded in the Geometric Priors described in Section~\ref{subsec:embedding_router}, leverages the rich cross-modal representations captured by pretrained facial cross-attention. Specifically, as shown in Figure \ref{fig3}, we utilize the query and key outputs from a pretrained multi-head face cross-attention module as inputs to the router module. Both query and key undergo layer-wise linear transformations followed by reshaping operations that preserve the head dimension. The attention weights, computed via matrix multiplication between transformed queries and keys, encode the correspondence between facial features and visual tokens. These weights are enhanced with 3D RoPE positional encoding~\cite{su2024roformer} before processing through multiple spatio-temporal attention layers that model interdependencies among visual tokens. A final linear layer with softmax activation generates the predicted mask output. This design holistically captures both isolated feature-token correspondences and spatio-temporal token dependencies, which are essential for producing precise, smooth masks.
\section{Teacher-Forcing Training Strategy}
We observed that training the router module jointly with the entire denoising process often leads to a trivial solution, where visual tokens become easy to classify but lack meaningful visual representation. Meanwhile, the router's training depends on the model being sufficiently adapted to the data distribution to perform high-quality denoising. Without accurate router predictions, the diffusion model progressively loses its ability to effectively integrate and utilize the provided conditions.

Thus,we employ a teacher-forcing training strategy, where we force the input of mask-guided cross attention in the transformer diffusion model to be the ground truth mask value, and detach the computational graph of the router module from the denoising process.
To enhance noise robustness, we apply a dropout operation and add Gaussian noise to the forced mask values, effectively serving as a form of data augmentation.